\section{Conclusion}\label{sec:conclusion}
The rung decides the verdict: RL beats OLLA, barely beats calibrated rules, and loses to a supervised one-step learner. Against OLLA, much of RL's gain came from calibration and convergence time. Against rules calibrated on its own inputs, a well-designed RL policy kept a small but robust edge. Against a supervised learner trained only on observed ACK/NACK, RL lost, with the same actions, across training draws and at two neighbouring operating points. It lost again on fresh traces evaluated after all policies were frozen. On full-buffer traffic with HARQ delay, RL added nothing over a calibrated offset.

For practitioners, the lesson is to calibrate the rules first and to try a one-step learner before deep RL, budgeting for its exploration and its decision time. For researchers, it is to report such comparators beside OLLA. The open question is where sequential learning starts to pay off: with deadlines that allow several retransmissions, with stronger coupling between users, or over the air.

\section*{Code, data and pre-registration availability}
\placeholder{Release DOI. The release will contain the simulator, the trace and CRN generators with all seeds, the per-seed and per-trace results of every scheme, the calibrated tables with their hashes, the case registry, the analysis scripts, and the frozen pre-registration texts with their SHA-256 hashes and freeze commits.}

\section*{Acknowledgments}
\placeholder{Funding sources and grant numbers; conflicts of interest; AI-assistance disclosure in the wording required by the journal.}
